% titel: Blatt 5, A1: Eigenschaften von Funktionen
% skript: Kap. 2 (Funktionen, Gerade/Ungerade, Monotonie, Umkehrfunktion)

% Teil I: Werkzeuge, kein Lösungsweg.

Eigenschaften von Funktionen (gerade/ungerade, Monotonie, Umkehrbarkeit)
analysieren und an Beispielen prüfen sowie einen Satz über die
Umkehrfunktion verstehen und begründen.

\begin{werkzeugkasten}[Gerade und ungerade Funktionen -- eine Spiegel-Frage]
Eine Funktion $f$ auf einer um $0$ symmetrischen Definitionsmenge heisst
\begin{itemize}
  \item \textbf{gerade}, wenn $f(x)=f(-x)$ für alle $x$ gilt. Der Graph ist
        eine Spiegelung seiner selbst an der $y$-Achse: Was rechts von $0$
        passiert, spiegelt sich exakt nach links.
  \item \textbf{ungerade}, wenn $f(-x)=-f(x)$ für alle $x$ gilt. Der Graph
        ist eine Spiegelung an der $y$-Achse \emph{und} zusätzlich an der
        $x$-Achse, zusammen also eine Drehung um $180^\circ$ um den
        Ursprung.
\end{itemize}
Der Name kommt von den typischen Monomen: $1,\,x^2,\,x^4$ sind gerade,
$x,\,x^3,\,x^5$ ungerade. Bei einem Produkt erbt man die Symmetrie als \emph{Produkt der Symmetrien}:
Gerade Monome hinterlassen beim Ersetzen von
$x$ durch $-x$ kein Minuszeichen, ungerade schon, also
$(\text{gerade})\cdot(\text{ungerade})=(\text{ungerade})$ und so weiter.
Man prüft aber immer \emph{direkt durch Einsetzen von $-x$}, nicht durch
Raten am Aussehen der Formel.

\textit{Skript:} Kap. 2 (Eigenschaften von Funktionen).
\end{werkzeugkasten}

\begin{werkzeugkasten}[Monotonie und Umkehrfunktion -- warum Injektivität reicht]
$f$ heisst \textbf{streng monoton steigend}, falls aus $x_1<x_2$ stets
$f(x_1)<f(x_2)$ folgt, \textbf{streng monoton fallend} bei
$f(x_1)>f(x_2)$. Anschaulich: Beim Durchlaufen der $x$-Achse bewegt sich
der Funktionswert nur in eine Richtung, es gibt kein Zurück. Genau das macht
$f$ \textbf{injektiv}: Jeder Funktionswert $y$ hat \emph{höchstens ein}
Urbild $x$.

Nur injektive Funktionen sind umkehrbar, denn die Umkehrfunktion $g$ fragt
\emph{Welches $x$ liefert diesen $y$-Wert?}. Hätte $y$ zwei Urbilder, wäre
die Antwort nicht eindeutig und $g$ keine Funktion. Dass die Umkehrfunktion
einer steigenden Funktion wieder steigt, folgt per \textbf{Kontraposition}:
Wäre $g$ irgendwo nicht steigend, so ergäbe das Anwenden der steigenden
Funktion $f$ einen Widerspruch zur Eindeutigkeit.

\textit{Skript:} Kap. 2 (Monotonie, Umkehrfunktion).
\end{werkzeugkasten}

\begin{achtung}[Typische Fallen]
\begin{itemize}
  \item \textbf{Definitionsbereich vergessen:} $f_1=\sqrt{1-x^2}$ ist nur
        auf $[-1,1]$ erklärt ($1-x^2\ge 0$). Symmetrie hat nur auf einer um
        $0$ symmetrischen Menge einen Sinn.
  \item \textbf{Nicht aussehen lassen:} Aus einem Quadratterm nicht allein
        auf gerade schliessen; die Wurzel reicht die Symmetrie zum Beispiel
        durch, ohne dass man am Vorzeichen rät.
  \item \textbf{Bijektivität und Monotonie verwechseln:} Bijektiv
        $[0,1]\to[0,1]$ erzwingt keine Monotonie ($x\mapsto 1-x$ ist
        bijektiv, aber fallend).
  \item \textbf{Beim Umkehrsatz die Kontraposition:} Man nimmt an, $g$ sei
        nicht steigend, und führt das zu einem Widerspruch -- statt die
        Richtung direkt zu \emph{erzwingen}.
\end{itemize}
\end{achtung}
